%% novelty_paragraph.tex
%% "What is new vs. prior work" — for the Visual RL paper
%% "Forecasting Competence Boundaries in Visual RL Without Rollouts"
%%
%% Covers: Lyu et al. 2024, kNN/Mahalanobis OOD, Majumdar et al. 2025,
%%          Romer et al. 2025 (FIPER), Farid et al. RSS 2022, Burns et al.,
%%          and affine-invariant Riemannian metrics (Moakher 2005).

\paragraph{Novelty relative to prior work.}
Several prior works motivate or relate to our approach, but each addresses a
strictly different problem.
%
\textbf{Lyu et al.~\cite{lyu2024generalization}} derive a theoretical upper bound on
the visual RL generalization gap as a function of representation drift, which
directly motivates our OOD severity score.
Critically, however, their bound is an \emph{upper bound that motivates our
score; it does not prove forecasting accuracy in practice.}
We show empirically where and when this bound fails to serve as a reliable
forecaster: on the smooth Visual Pursuit domain it correlates strongly with
performance ($r = -0.8853$), but on the discrete Procgen CoinRun benchmark
the same metric collapses ($\rho = -0.1143$, $p = 0.45$), confirming that the
theoretical bound does not translate automatically into a practical forecasting
tool.
%
Standard \textbf{kNN and Mahalanobis distance OOD scores}~\cite{lee2018simple}
are designed for classification: they detect whether an input is anomalous
with respect to class-conditional feature distributions.
We instead apply distance-based scores to the \emph{policy feature space} of a
visual RL agent, pair them with subspace rank selection via Grassmann geodesics,
and use them to \emph{forecast} the severity of performance degradation across
a continuous perturbation axis — not merely to flag distributional anomalies.
%
\textbf{Majumdar et al.~\cite{majumdar2025predictive}} (RoboART / Predictive Red
Teaming) also predict policy failure without physical hardware rollouts, but
their pipeline relies on generative image editing of \emph{observed rollouts
from the policy} to construct off-nominal test conditions.
Our approach requires no rollout access whatsoever: we operate entirely on
single static observations and pre-computed feature statistics, making it
applicable in deployment scenarios where policy execution is infeasible or
prohibited.
%
\textbf{Römer et al.~\cite{romer2025fiper}} (FIPER) address runtime failure
prediction for generative imitation learning policies, targeting binary
detection of impending task failure.
Our target is different: we forecast the \emph{quantitative severity} of the
competence boundary (the expected drop in tracking performance as a function
of shift intensity), not whether a single execution will succeed or fail.
%
\textbf{Farid et al.~\cite{farid2022failure}} provide PAC-Bayes statistical
guarantees for failure prediction given a \emph{history} of sensor observations
from a running policy.
They define safety regions in observation space; we instead define competence
regions in feature-drift space and forecast their boundaries without requiring
online policy execution or a labelled failure dataset.
%
\textbf{Burns et al.~\cite{burns2023weak}} study generalization in a
weak-supervision setting fundamentally different from ours; their framework
is not directly applicable to online-free competence boundary estimation in
visual RL~\textsuperscript{$\dagger$}.
%% ^\dagger NOTE TO AUTHORS: Confirm which Burns et al. paper is intended
%%   and update this sentence accordingly. If a different Burns et al. is
%%   cited, revise the contrast accordingly.
%
Finally, we apply the \textbf{affine-invariant Riemannian metric}
(AIRM)~\cite{moakher2005differential} to measure geodesic distance between
policy feature covariance matrices, treating them as points on the symmetric
positive-definite (SPD) manifold.
Unlike Euclidean covariance distance, AIRM is invariant to affine
reparametrizations of the feature space, making it a principled and
scale-equivariant measure of representation drift.
To our knowledge, this is the first application of AIRM to policy feature
covariance matrices for the purpose of OOD distance estimation in visual RL.
